\documentclass[../main.tex]{subfiles}

\begin{document}

\section{Resultados obtenidos}

Los resultados presentados a continuación responden a las cuatro hipótesis planteadas en la Sección~3 y se organizan en cinco bloques: la descripción general del experimento (§5.1), el poder predictivo de QUILL (§5.2, H1), el predominio del Poder Discriminativo (§5.3, H2), las propiedades lingüísticas de las pistas (§5.4, H3) y el efecto de la modalidad experimental (§5.5, H4).

En este capítulo se presentan los resultados empíricos derivados de la implementación y despliegue del sistema de generación de pistas. El análisis se divide en dos fases consecutivas, comenzando con una evaluación intrínseca de la calidad de las pistas generadas basada en las métricas de calidad semántica, definidas en el capítulo anterior, y una validación extrínseca basada en el comportamiento real de los usuarios para poder inferir en el peso que tendrá cada una de estas en la fricción cognitiva y su utilidad didáctica o discriminativa hacia el usuario. Seguido a esto, realizaremos el estudio de las cualidades sintácticas de la oración, esperando obtener datos en cuanto a la relación entre sintáctica y semántica de la frase.

Para la validación experimental, contamos con un conjunto de datos compuesto por un total de $N=3412$ interacciones de usuarios únicos. Este volumen se divide metodológicamente en dos etapas: una primera recolección de $N_{abierto}=2293$ interacciones en un entorno de evocación libre, y una segunda de $N_{cerrado}=1119$ interacciones en un entorno de reconocimiento deductivo. Este conjunto de datos ($D_{logs}$) permite contrastar la eficacia teórica de las métricas propuestas ($Q$, PD, AMB, etc.) frente a la tasa de éxito real (\textit{Ground Truth}) en un entorno de juego no controlado.

El objetivo principal de esta sección es inferir en qué hace a una pista, como unidad de texto, de calidad o ``util'' en términos de usuario; pasando primero por las etapas de validación para las que hemos recopilado información útil de nuestros usuarios

\subsection{Descripción general del experimento}

Antes de evaluar la interacción con el usuario, es fundamental analizar las propiedades estadísticas de las pistas generadas por el modelo. Para ello, se calcularon las métricas vectoriales y formales sobre la totalidad del corpus de pistas disponibles en el sistema ($N_{pistas} = 360$). Proponemos una jerarquización teórica basada en la idoneidad pedagógica y la capacidad informativa, remarcando que esto se ancla en la subjetividad del criterio profesional. El siguiente orden decreciente refleja la importancia estimada de cada atributo para la generación de pistas efectivas

\begin{itemize}
    \item \textbf{Ganancia de información:} estimada como la métrica prioritaria. Una pista no debe limitarse a señalar un resultado, sino que debe aportar una carga semántica significativa. El objetivo es reducir la entropía lo suficiente para orientar al usuario, pero manteniendo un margen de incertidumbre que estimule el proceso cognitivo y el aprendizaje, evitando dar la solución de forma inmediata, por esto, se le dota de un 30\% del peso en la asignación de calidad.

    \item \textbf{Poder Discriminativo:} pensado fundamental para la precisión, la métrica evalúa la capacidad de la pista para aislar el concepto objetivo y diferenciarlo de otros elementos semánticamente cercanos o ``distractores''. Una alta discriminación asegura que la descripción no sea aplicable a múltiples candidatos erróneos. Por este motivo, el 20\% de la calidad responde a este indicador

    \item \textbf{Ambigüedad Relativa:} cuya minimización es crítica, ya que una pista excesivamente vaga, o que proporcione datos genéricos compartidos por un amplio conjunto de objetivos, carece de valor educativo. La ambigüedad no solo confunde, sino que diluye el propósito de la pista al no permitir descartar opciones incorrectas con certeza. Así, es responsable de otro 20\% de la calidad.

    \item \textbf{Consistencia:} se refiere a la coherencia interna y la precisión temática. Una pista que intenta abarcar demasiados aspectos corre el riesgo de perder el foco, generando ruido informativo, o de volverse contradictoria. La consistencia garantiza que la información proporcionada sea fiable y pertinente al contexto específico del objetivo. Se le otorga un 15\% de la responsabilidad de la asignación de puntaje.

    \item \textbf{Lecturabilidad:} una redacción clara y sintácticamente sencilla reduce la carga cognitiva necesaria para decodificar el texto. Al facilitar la comprensión lectora, permitimos que el usuario dedique la mayor parte de sus recursos mentales al proceso de aprendizaje y resolución del problema, en lugar de a la interpretación de la sintaxis. Al no ser tan importante en la dispensación de información, obtiene un 15\% del potencial.
\end{itemize}

\subsubsection{Distribución de las Métricas de Calidad}

La Figura \ref{fig:distribucion_metricas} muestra los histogramas de frecuencia para cada una de las componentes de la función de evaluación, así como para el Score Final $Q$.

% TODO: Inserta aquí el nombre de tu archivo con los 7 histogramas
\begin{figure}[H]
    \centering
    \includegraphics[width=1.0\textwidth]{imagenes/histogramas_metricas.png} 
    \caption{Distribución de frecuencias de las métricas individuales y el Score Q.}
    \label{fig:distribucion_metricas}
\end{figure}

Del análisis visual de estas distribuciones se extraen tres conclusiones diagnósticas sobre el comportamiento del generador:

\begin{itemize}
    \item \textbf{Normalidad del Score $Q$:} $Q$ es la función de calidad compuesta QUILL —definida en la Sección~4.5— que integra linealmente las cinco métricas semánticas (PD, AMB, nGAIN, CONS, LEC) ponderadas y protegidas por la restricción CLAMP. La distribución de la puntuación final $Q$ (gráfica inferior derecha) presenta una forma de campana de Gauss casi perfecta, centrada en el intervalo $[0.3, 0.4]$. La distribución aproximadamente gaussiana de Q sugiere que la métrica proporciona una variabilidad suficiente y evita concentraciones en los extremos. No obstante, esta propiedad por sí sola no garantiza su validez, que debe evaluarse en función de su capacidad para explicar o predecir el comportamiento del usuario.
    
    \item \textbf{Poder Discriminativo (PD):} La distribución de PD presenta una ligera asimetría positiva, con valores que oscilan entre $-0.16$ y $+0.30$ y una mediana próxima a $+0.05$. Aunque la mayoría de las pistas generadas exhiben un PD positivo —indicando que, en promedio, se acercan más al objetivo que a los distractores—, aproximadamente el 20\,\% del corpus presenta PD negativo, lo que revela casos en los que el contenido semántico del texto no es suficientemente diferenciador dentro del espacio de candidatos.

    \item \textbf{Ganancia de Información Normalizada (nGAIN):} El histograma de nGAIN muestra una distribución unimodal centrada en torno a 0.4, indicando que las pistas generadas reducen, en promedio, aproximadamente el 40\,\% de la entropía inicial del espacio de búsqueda. Los valores extremos cercanos a 0 corresponden a pistas que prácticamente no concentran la distribución de probabilidad sobre ningún candidato, mientras que los valores próximos a 0.8 se asocian a pistas que proyectan la mayor parte de la masa de probabilidad sobre un subconjunto muy reducido de alternativas.

    \item \textbf{Sesgo de Ambigüedad (AMB):} Se observa que la mediana de la ambigüedad (AMB) se sitúa en torno a 0,60, lo que sugiere que muchas pistas presentan similitud semántica con múltiples candidatos. Este comportamiento indica que, en este contexto, resulta frecuente que las pistas no distingan claramente al objetivo dentro del conjunto de posibles respuestas. En este sentido, la ambigüedad emerge como un factor relevante a controlar, ya que niveles elevados pueden dificultar el proceso de discriminación por parte del usuario.

    \item \textbf{Saturación de Métricas Formales:} Tanto la consistencia (CONS) como la lecturabilidad (LEC) presentan distribuciones sesgadas hacia valores altos (mediana de CONS $\approx$ 0,75 y LEC cercana a 1,0), lo que indica que la mayoría de las pistas generadas cumplen adecuadamente criterios básicos de corrección sintáctica y facilidad de lectura. Este comportamiento sugiere que estas métricas presentan una baja variabilidad en el conjunto analizado, por lo que su capacidad para discriminar entre pistas de distinta calidad es limitada. En este sentido, CONS y LEC parecen desempeñar un papel más cercano a indicadores de adecuación mínima que a factores determinantes en la efectividad de las pistas.

    \item \textbf{Latencia de CLAMP}: La condición de parada CLAMP no se activa en los datos analizados, lo que indica que el modelo de lenguaje ha respetado consistentemente las restricciones impuestas en el proceso de generación, evitando la inclusión explícita de la entidad objetivo en las pistas. Este comportamiento sugiere que CLAMP actúa principalmente como un mecanismo de control que garantiza la validez de las pistas generadas, más que como un factor discriminante en su calidad.
\end{itemize}

\subsubsection{Análisis de Correlación e Independencia de Métricas}

Para validar la ortogonalidad de las métricas propuestas y detectar posibles redundancias, se calculó la matriz de correlación de Pearson entre todas las variables del sistema (Ver Figura \ref{fig:heatmap_correlacion}).

% TODO: Inserta aquí el nombre de tu archivo con el Heatmap rojo/azul
\begin{figure}[H]
    \centering
    \includegraphics[width=0.8\textwidth]{imagenes/heatmap_correlacion.png}
    \caption{Matriz de correlación de Pearson entre las métricas de evaluación.}
    \label{fig:heatmap_correlacion}
\end{figure}

El mapa de calor revela relaciones críticas para el diseño de la función de agregación:

\begin{enumerate}
    \item \textbf{Validación de $Q$ ($r=0.92$):} Existe una correlación positiva muy fuerte entre el Poder Discriminativo (PD) y el Score Final ($Q$), y una correlación negativa fuerte con la Ambigüedad ($-0.87$). Esto confirma que la fórmula $Q$ está alineada correctamente con los objetivos semánticos del sistema: premiar la especificidad y castigar la confusión.
    
    \item \textbf{Problema de Colinealidad ($r=-0.80$):} Se detectó una correlación inversa significativa de $-0.80$ entre PD y AMB. Esto indica una redundancia en la evaluación: las pistas que se parecen mucho al objetivo tienden a parecerse poco al rival. Esta colinealidad sugiere que asignar pesos elevados a ambas métricas simultáneamente podría estar ``contando dos veces'' el mismo fenómeno semántico, inflando artificialmente la varianza del score.
    
    \item \textbf{Independencia de la Ganancia (nGAIN):} La correlación entre la Ganancia de Información y la Ambigüedad es moderada ($-0.29$). Esto es un hallazgo positivo, pues demuestra que nGAIN captura información única sobre la reducción de entropía global que no está siendo medida por la comparación directa con el rival (AMB).

\end{enumerate}

A pesar de la colinealidad detectada, ambas métricas capturan aspectos complementarios y se mantienen en la función QUILL. PD mide la brecha media entre la pista y el conjunto global de distractores; AMB se centra en el rival semántico más cercano —el caso peor—. Una pista puede presentar un PD aceptable y, sin embargo, una AMB crítica si un único distractor está muy próximo al objetivo en el espacio de embeddings, situación que PD no detecta. En la función QUILL, el término $(1 - 1{,}5\,\text{AMB})$ penaliza específicamente esa proximidad con el rival más cercano, aportando una dimensión ortogonal a la media global capturada por PD.

Estos hallazgos descriptivos plantean la necesidad de ajustar los pesos de la fórmula $Q$ basándose no en la intuición, sino en la capacidad real de cada métrica para predecir el éxito del usuario, proceso que se detalla en la siguiente sección.

\subsubsection{Validación empírica de la función CLAMP}

Un hallazgo metodológico relevante emergió al analizar la distribución empírica de la similitud directa entre cada pista y su objetivo (PROX). La formulación original del CLAMP establecía un umbral arbitrario de $\delta = 0.90$: únicamente las pistas con $\text{PROX} > 0.90$ serían penalizadas. Sin embargo, el valor máximo de PROX observado en el corpus completo de 360 pistas es $0.678$, lo que implica que \textbf{ninguna pista superó jamás el umbral}: el CLAMP original devolvía 1.0 para la totalidad del corpus, actuando como un multiplicador neutro sin efecto práctico sobre ninguna de las métricas ni sobre la selección de pistas.

Este hallazgo motivó la reformulación del CLAMP descrita en el Capítulo~4, calibrada automáticamente sobre la distribución empírica de PROX mediante una función sigmoide centrada en dos sigmas. La Figura~\ref{fig:clamp_calibration} muestra la distribución empírica de PROX (izquierda) y la curva CLAMP resultante (derecha).

\begin{figure}[H]
    \centering
    \includegraphics[width=0.92\textwidth]{imagenes/clamp_calibration.png}
    \caption{Izquierda: histograma de PROX en el corpus de 360 pistas, con $\mu = 0.273$ y el umbral $\mu + 2\sigma = 0.579$ marcados. Derecha: función CLAMP sigmoide calibrada sobre el rango empírico, con los 360 puntos del corpus superpuestos.}
    \label{fig:clamp_calibration}
\end{figure}

La nueva formulación asigna factores CLAMP en el intervalo $[0.34, 0.99]$, con media $\bar{C} = 0.843$. La mayoría de las pistas reciben un descuento moderado, proporcional a qué tan por encima de la media del corpus se encuentra su similitud directa al objetivo. Las pistas con PROX excepcionalmente alto quedan penalizadas de forma continua y proporcional, sin la discontinuidad abrupta del umbral manual original.

\subsection{H1 --- Poder predictivo de QUILL}

Para evaluar la utilidad funcional real de las métricas definidas, se utilizó el conjunto de interacciones $D_{logs}$ de la primera fase experimental ($N = 2293$), segmentando los registros en:

\begin{itemize}
\item Clase positiva ($y=1$): aciertos directos.
\item Clase negativa ($y=0$): fallos, solicitudes adicionales o abandono.
\end{itemize}

Dado el leve desbalance del conjunto (aproximadamente 1:1.5), se empleó una regresión logística con ponderación de clases (\texttt{class\_weight='balanced'}). Todas las variables independientes fueron previamente estandarizadas mediante \texttt{StandardScaler} (transformación z-score), permitiendo comparar directamente las magnitudes relativas de los coeficientes estimados.

Es importante destacar que el objetivo del modelo no es construir un predictor conductual de alta precisión, sino estimar la contribución marginal relativa de cada métrica a la probabilidad de acierto.

\subsubsection{Validación Cruzada y Capacidad Discriminativa}

Para evitar dependencia de un único particionado train/test, se realizó validación cruzada estratificada 5-fold.

Los resultados medios obtenidos fueron:

\begin{center}
\begin{tabular}{|l|c|}
\hline
Accuracy & $0.546 \pm 0.026$ \\
AUC-ROC & $0.565 \pm 0.027$ \\
\hline
\end{tabular}
\end{center}

Estos valores indican:

\begin{itemize}
\item Capacidad discriminativa modesta pero consistente ($\text{AUC} > 0.5$).
\item Baja varianza entre folds ($\text{std} < 0.03$), lo que sugiere estabilidad estructural del modelo.
\end{itemize}

Aunque la mejora sobre el azar es limitada, este resultado es coherente con la naturaleza del problema: la probabilidad de acierto humano depende de múltiples factores no observados (conocimiento previo, experiencia, estrategia), mientras que el modelo solo utiliza características de la pista.

\begin{figure}[H]
    \centering
    \includegraphics[width=0.75\textwidth]{imagenes/roc_curve_cv.png}
    \caption{Curva ROC media en validación cruzada 5-fold con banda de desviación estándar.}
    \label{fig:roc_cv}
\end{figure}

La curva se sitúa consistentemente por encima de la diagonal, confirmando señal predictiva real aunque de magnitud moderada.

\subsubsection{Bondad de Ajuste Global}

El pseudo-$R^2$ de McFadden obtenido fue:

\[
R^2_{McFadden} = 0.0102
\]

Este valor es bajo en términos absolutos, indicando que el modelo explica una fracción reducida de la variabilidad total frente al modelo nulo. Sin embargo, este resultado es esperable dado que únicamente se modelan propiedades del estímulo y no variables del jugador.

Para contrastar formalmente la relevancia global del modelo se realizó un test de razón de verosimilitud:

\[
G^2 = 30.10, \quad p = 1.41 \times 10^{-5}
\]

El resultado significativo confirma que, colectivamente, las métricas contienen información predictiva real sobre la probabilidad de acierto.

Adicionalmente, se calculó el Brier Score:

\[
BS = 0.245
\]

Dado que un modelo puramente aleatorio obtiene aproximadamente $0.25$, el valor obtenido ($0.245 < 0.25$) indica que el modelo mejora ligeramente al azar en términos de calibración probabilística, resultado coherente con la capacidad discriminativa modesta observada en el AUC.

Para ilustrar de manera concreta el poder predictivo de QUILL, se presentan las tres pistas con mayor puntuación $Q$ inicial del corpus:

\begin{table}[H]
\centering
\small
\begin{tabular}{llcccccc}
\hline
\textbf{Persona} & \textbf{Q} & \textbf{PD} & \textbf{AMB} & \textbf{nGAIN} & \textbf{CONS} & \textbf{LEC} & \textbf{NER} \\
\hline
Buda Gautama     & 0.507 & 0.285 & 0.360 & 0.638 & 0.756 & 0.995 & 2 \\
Napoleón Bonaparte & 0.489 & 0.304 & 0.410 & 0.691 & 0.736 & 0.996 & 3 \\
Winston Churchill  & 0.478 & 0.287 & 0.441 & 0.546 & 0.745 & 0.995 & 1 \\
\hline
\end{tabular}
\caption{Pistas con mayor puntuación QUILL del corpus (pesos iniciales, $N=360$).}
\label{tab:top_hints}
\end{table}

\textbf{Buda Gautama} ($Q = 0.507$): \textit{``Nació en Lumbini en el seno de una familia aristocrática en la república Sakia.''} Esta pista encarna el perfil óptimo definido por el análisis OLS: dos entidades geográficas altamente específicas (Lumbini, Sakia), un único verbo conjugado y ausencia de adjetivos valorativos. Las entidades actúan como anclajes factuales inequívocos que recortan de forma drástica el espacio de candidatos (nGAIN = 0.638), mientras la baja carga verbal y atributiva garantiza la máxima lecturabilidad (LEC = 0.995).

\textbf{Napoleón Bonaparte} ($Q = 0.489$): \textit{``Nació en la isla de Córcega, poco después de su anexión por Francia.''} Ejemplo paradigmático del principio de especificidad nominal: densidad de adjetivos nula (ADJ\_Ratio = 0.000) y tres entidades nombradas. La combinación de dato geográfico con evento histórico concreto produce un poder discriminativo elevado (PD = 0.304) sin sacrificar la accesibilidad cognitiva.

\textbf{Winston Churchill} ($Q = 0.478$): \textit{``Se unió al ejército británico en 1895 y vio acción en la India británica.''} La presencia de un anclaje temporal específico (1895) y una entidad geográfica (India) produce un hecho difícilmente generalizable a otros candidatos. La consistencia interna es alta (CONS = 0.745): ambas cláusulas orbitan sobre el mismo eje temático del servicio militar.

Los tres casos comparten las mismas propiedades estructurales: baja densidad de verbos (VERB\_Ratio $< 0.15$), ausencia o mínima presencia de adjetivos, y al menos una entidad geográfica o temporal que actúa como restricción lógica directa en el espacio semántico.

\subsection{H2 --- Predominio del Poder Discriminativo}

\subsubsection{Coeficientes del Modelo}

La Tabla \ref{tab:coeficientes_logistica} muestra los coeficientes estandarizados estimados.

\begin{table}[H]
\centering
\begin{tabular}{|l|c|c|}
\hline
\textbf{Métrica} & \textbf{Coeficiente ($\beta$)} & \textbf{Signo} \\ \hline
PD & +0.110 & Positivo \\
LEC & -0.070 & Negativo \\
AMB & -0.080 & Negativo \\
nGAIN & +0.056 & Positivo \\
CONS & -0.008 & Nulo \\
\hline
\end{tabular}
\caption{Coeficientes estandarizados de la regresión logística ($N=2293$).}
\label{tab:coeficientes_logistica}
\end{table}

Cabe señalar que, aunque el modelo global resulta significativo, los coeficientes individuales no alcanzan significancia estadística estricta ($p<0.05$) cuando se estiman mediante \textit{statsmodels}. Esto es consistente con la presencia de colinealidad parcial entre métricas y con la dificultad intrínseca del problema.

En consecuencia, la interpretación debe realizarse en términos de contribución relativa y no como evidencia causal individual.

\subsubsection{Interpretación Empírica}

\begin{enumerate}

\item \textbf{Predominio del poder discriminativo:} PD presenta el mayor coeficiente positivo. Esto confirma que la separación efectiva respecto al competidor más cercano es el principal determinante observable del éxito.

\item \textbf{Efecto inverso de la lecturabilidad:} El coeficiente negativo de LEC no implica que la complejidad lingüística sea deseable per se, sino que la métrica Flesch penaliza vocabulario técnico y específico, el cual resulta útil para discriminar.

\item \textbf{Penalización por ambigüedad:} AMB mantiene efecto negativo consistente, reforzando la necesidad de minimizar similitudes estructurales con candidatos rivales.

\item \textbf{Aporte positivo de nGAIN:} La reducción global de entropía contribuye al éxito, aunque en menor magnitud que PD.

\item \textbf{Rol marginal de CONS:} La baja varianza observada en coherencia interna limita su capacidad discriminativa en este conjunto concreto.

\end{enumerate}

\subsubsection{Reestimación de Pesos}

Con el objetivo de alinear la función de evaluación con la evidencia empírica, se recalcularon los pesos normalizando la magnitud absoluta de los coeficientes:

\[
w_i = \frac{|\beta_i|}{\sum |\beta_k|}
\]

Esta estrategia permite construir un índice compuesto basado en contribución relativa observada, aunque no implica inferencia causal individual fuerte.

La nueva formulación queda definida como:

\[
Q_{opt}(h) =
\text{CLAMP}(h) \cdot
\left[
0.335\,\text{PD}
+ 0.215\,(1 - 1.5\,\text{AMB})
+ 0.139\,\text{nGAIN}
+ 0.009\,\text{CONS}
+ 0.303\,\text{LEC}
\right]
\]

\begin{figure}[H]
    \centering
    \includegraphics[width=0.7\textwidth]{imagenes/q_comparison_histogramas.png}
    \caption{Histogramas superpuestos de $Q$ Inicial (pesos teóricos) y $Q$ Mundo Abierto (pesos empíricos derivados de la Fase 1). La ligera diferencia en las distribuciones refleja el rebalanceo de pesos inducido por la regresión logística.}
    \label{fig:comparativa_pesos}
\end{figure}

En conjunto, los resultados muestran que la calidad efectiva de una pista depende principalmente de su capacidad para aislar el objetivo respecto a competidores cercanos, manteniendo suficiente especificidad léxica para evitar generalidad excesiva.

Para ilustrar el impacto del bajo poder discriminativo, se presentan las tres pistas con menor puntuación $Q$ inicial y las reformulaciones propuestas:

\begin{table}[H]
\centering
\small
\begin{tabular}{llcccccc}
\hline
\textbf{Persona} & \textbf{Q} & \textbf{PD} & \textbf{AMB} & \textbf{nGAIN} & \textbf{CONS} & \textbf{LEC} & \textbf{NER} \\
\hline
Molière (pista 1)  & -0.068 & -0.144 & 1.910 & 0.232 & 0.758 & 0.996 & 2 \\
Molière (pista 2)  &  0.132 & -0.155 & 1.190 & 0.211 & 0.744 & 0.983 & 0 \\
Marie Curie        &  0.183 & -0.060 & 0.875 & 0.362 & 0.738 & 0.293 & 0 \\
\hline
\end{tabular}
\caption{Pistas con menor puntuación QUILL del corpus (pesos iniciales, $N=360$).}
\label{tab:bottom_hints}
\end{table}

\textbf{Molière (pista 1)} ($Q = -0.068$): \textit{``Nació en una familia próspera y estudió en el Collège de Clermont.''} El PD negativo ($-0.144$) indica que el embedding de la frase es semánticamente más cercano a otros candidatos del corpus que a Molière. El descriptor ``familia próspera'' es aplicable a decenas de figuras históricas y añade ruido entrópico puro. La AMB = 1.910 es la mayor del corpus: el mejor candidato rival atrae casi el doble de señal semántica que el objetivo real. El fallo es estructural: la pista aporta información veraz pero no discriminativa dentro del espacio de candidatos.

\textbf{Reformulación propuesta}: \textit{``Dramaturgo francés del siglo XVII, autor de \textit{El avaro}, \textit{El enfermo imaginario} y \textit{Tartufo}.''} Esta versión elimina los verbos de estado genéricos y sustituye la descripción biográfica difusa por tres entidades nombradas concretas (títulos de obra propios del objetivo), elevando NER de 2 a 3 y convirtiendo la pista en una restricción lógica directa.

\textbf{Molière (pista 2)} ($Q = 0.132$): \textit{``Su obra \textit{Tartufo} fue prohibida por el Parlamento francés debido a su crítica a la hipocresía religiosa.''} El PD vuelve a ser negativo ($-0.155$): el campo semántico de censura religiosa y parlamento activa representaciones compartidas con numerosos personajes del siglo XVII europeo. A pesar de referirse a un hecho específico de Molière, la pista no se ancla semánticamente al objetivo porque el LLM ha producido una formulación centrada en el contexto histórico-político y no en la figura del autor. La ausencia completa de entidades nombradas en el lado del autor (NER = 0) evidencia el patrón de dispersión temática identificado por los modelos OLS.

\textbf{Reformulación propuesta}: \textit{``Falleció en 1673 durante la representación de su última obra en el Palais-Royal de París.''} Dos entidades específicas (1673, Palais-Royal de París) anclan la pista a un evento único e irrepetible, corrigiendo el PD negativo al proyectar la pista hacia el espacio semántico propio del objetivo.

\textbf{Marie Curie} ($Q = 0.183$): \textit{``Desarrolló técnicas para el aislamiento de isótopos radiactivos.''} El fallo principal no es la precisión factual, sino la lecturabilidad extremadamente baja (LEC = 0.293), consecuencia directa de la acumulación de polisílabos de alta complejidad silábica. El índice de Flesch-Szigriszt penaliza severamente esta densidad, independientemente de la precisión del contenido. El PD negativo ($-0.060$) adicional indica que términos como ``isótopos radiactivos'' activan campos semánticos de física nuclear compartidos con otros científicos del corpus.

\textbf{Pista alternativa propuesta}: \textit{``Primera mujer en ganar el Premio Nobel, en 1903, junto a Pierre Curie.''} La nueva versión reemplaza el vocabulario técnico polisílabo por una formulación basada en dos entidades nombradas accesibles (Premio Nobel, Pierre Curie) y una fecha precisa.

\subsection{H3 --- Propiedades lingüísticas y calidad semántica}

A partir de los resultados empíricos obtenidos en la sección anterior, es posible formalizar con mayor precisión el concepto de pista como unidad funcional dentro del sistema. Definimos una pista como una unidad composicional de información expresada en lenguaje natural, condicionada por un objetivo latente $t$, cuya función es modificar el estado epistemológico del receptor reduciendo la entropía del espacio de estados posibles $\mathcal{S}$ sin revelar explícitamente dicho objetivo.

Sea $\mathcal{S} = \{s_1, \dots, s_N\}$ el conjunto de candidatos posibles y $P(s)$ la distribución inicial asumida por el jugador, que en ausencia de información se aproxima a una distribución uniforme. Una pista $h$ actúa como operador informacional que induce una actualización bayesiana:

\[
P(s \mid h) \propto P(h \mid s)\,P(s)
\]

De este modo, la pista puede interpretarse como un \textbf{filtro bayesiano lingüístico}: transforma una distribución \textit{a priori} relativamente plana en una distribución \textit{a posteriori} donde la masa de probabilidad se concentra progresivamente sobre el objetivo $t$. La calidad de la pista no reside únicamente en la magnitud de esta concentración, sino en la forma en que se produce dicha redistribución.

Esta reformulación mejora el baseline inicial (``unidad de información basada en lenguaje natural cuyo cometido es eliminar entropía'') en tres dimensiones fundamentales:

\begin{enumerate}

\item \textbf{Composicionalidad estructural.}
Una pista no es un dato atómico, sino una estructura lingüística compuesta por múltiples componentes (entidades nominales, propiedades atributivas, relaciones sintácticas y dependencias semánticas). Cada uno de estos elementos contribuye diferencialmente a la reducción de incertidumbre. Las regresiones OLS que se presentan a continuación permiten cuantificar qué características formales están asociadas con mayor poder discriminativo o mayor ganancia de información.

\item \textbf{Condicionalidad relacional.}
La calidad de una pista no es intrínseca al texto en sí mismo, sino relacional. Una misma secuencia lingüística puede ser altamente informativa para un objetivo $t_1$ y prácticamente irrelevante para $t_2$. La evaluación debe entenderse como una función ternaria texto–objetivo–campo semántico. Formalmente, $Q(h)$ depende de la estructura del espacio $\mathcal{S}$ y de la posición relativa del objetivo dentro de dicho espacio.

\item \textbf{No-trivialidad óptima.}
El objetivo de una pista no es colapsar la entropía a cero (lo cual equivaldría a revelar la respuesta), sino producir una reducción calibrada que mantenga un gradiente informativo. Una pista óptima equilibra discriminación y ambigüedad residual: reduce significativamente el conjunto plausible de estados sin eliminar completamente la incertidumbre. Esta tensión entre claridad y no-revelación constituye el núcleo del problema de optimización.

\end{enumerate}

Desde esta perspectiva, la pista puede descomponerse en dos niveles analíticos complementarios:

\begin{itemize}
\item Nivel \textbf{semántico-funcional}: mide el efecto global sobre la distribución $P(s)$ (PD, AMB, nGAIN).
\item Nivel \textbf{estructural-lingüístico}: caracteriza la forma en que dicha información es codificada (longitud, complejidad sintáctica, densidad nominal, subjetividad, etc.).
\end{itemize}

Mientras que la regresión logística permitió estimar qué métricas semánticas se asocian con el éxito observable del jugador, el análisis OLS que se presenta a continuación busca responder una pregunta distinta pero complementaria:

\begin{quote}
¿Qué propiedades formales del texto explican variaciones sistemáticas en las métricas semánticas de calidad?
\end{quote}

En otras palabras, si la regresión logística modela la \textit{utilidad externa} de la pista, la regresión OLS modela su \textit{anatomía interna}. Este análisis permite entender no solo qué métricas son relevantes, sino cómo emergen estructuralmente a partir de la configuración lingüística del enunciado.

Sea $X \in \mathbb{R}^{n \times p}$ la matriz de características lingüísticas estandarizadas y $Y_k$ cada una de las métricas de calidad ($k \in \{PD, AMB, nGAIN, LEC, CONS\}$). Para cada métrica se estima el modelo lineal:

\[
Y_k = X\beta_k + \varepsilon_k
\]

donde:

\begin{itemize}
\item $\beta_k$ representa el vector de coeficientes estandarizados,
\item $\varepsilon_k$ es el término de error,
\item se utilizan errores robustos HC3 para garantizar inferencia válida ante heterocedasticidad.
\end{itemize}

Todas las variables se transforman mediante estandarización Z-score, lo que permite interpretar cada coeficiente $\beta$ como el cambio en desviaciones estándar de la métrica ante un incremento de una desviación estándar en el predictor.

Antes de examinar los modelos individuales, es importante establecer una expectativa teórica:

\begin{itemize}
\item Las métricas puramente estructurales (como LEC) deberían mostrar mayor varianza explicada.
\item Las métricas semánticas profundas (PD, nGAIN) deberían presentar $R^2$ bajos pero potencialmente significativos.
\item Si AMB depende realmente del campo semántico relacional, debería resultar inexplicable desde variables superficiales.
\end{itemize}

Esta hipótesis constituye una prueba indirecta de la validez conceptual de las métricas diseñadas.

Las características independientes se agrupan en tres estratos conceptuales:

\begin{itemize}
\item \textbf{Anclaje nominal}: PROPN\_Ratio, NER\_Count, NUM\_Ratio, Max\_IDF.
\item \textbf{Carga narrativa}: VERB\_Ratio, ADJ\_Ratio.
\item \textbf{Estructura sintáctica}: Tree\_Depth.
\end{itemize}

Esta clasificación permitirá interpretar los resultados como dinámicas estructurales entre anclaje factual y dispersión narrativa.

A continuación se presentan cinco subsecciones independientes, cada una dedicada a una métrica específica. En cada caso se analizarán:

\begin{itemize}
\item Bondad de ajuste global (R², F-test).
\item Coeficientes estandarizados y significancia.
\item Diagnósticos (Q-Q plot, residuos vs ajustados).
\item Implicaciones estructurales derivadas de los coeficientes.
\end{itemize}

\subsubsection{Poder Discriminativo (PD)}

El Poder Discriminativo (PD) se define como:

\[
PD = \text{sim}(h,t) - \text{mean}(\text{sim}(h, s \neq t))
\]

donde la similitud se calcula mediante coseno en el espacio de embeddings. La métrica captura la separación relativa entre el objetivo y el conjunto de competidores activados por la pista.

\paragraph{Bondad de ajuste}

El modelo OLS estimado sobre $N = 360$ observaciones produce los siguientes resultados globales:

\begin{itemize}
\item $R^2 = 0.052$
\item $R^2$ ajustado = 0.033
\item $F = 3.458$
\item $p$-valor global = 0.00135
\end{itemize}

El modelo resulta globalmente significativo, aunque con una capacidad explicativa reducida (aproximadamente 5.2\% de la varianza). Esto indica que las características lingüísticas superficiales contienen señal estadística, pero limitada, en la explicación del PD.

\paragraph{Coeficientes estandarizados}

La Tabla siguiente resume los coeficientes estimados con errores robustos HC3:

\begin{table}[H]
\centering
\begin{tabular}{lccc}
\hline
Variable & $\beta$ & $p$-valor & Signif. \\
\hline
PROPN\_Ratio & 0.0248 & 0.655 & -- \\
NUM\_Ratio & 0.0006 & 0.993 & -- \\
Max\_IDF & -0.1000 & 0.010 & * \\
ADJ\_Ratio & 0.0991 & 0.071 & $\dagger$ \\
VERB\_Ratio & -0.1329 & 0.020 & * \\
Tree\_Depth & 0.0179 & 0.708 & -- \\
NER\_Count & 0.1082 & 0.068 & $\dagger$ \\
\hline
\end{tabular}
\caption{Coeficientes estandarizados del modelo OLS para PD (errores robustos HC3).}
\end{table}

(* $p<0.05$, $\dagger$ $p<0.10$)

La Figura \ref{fig:coef_pd} muestra los coeficientes estandarizados junto con sus intervalos de confianza.

\begin{figure}[H]
    \centering
    \includegraphics[width=0.9\textwidth]{imagenes/ols_PD_individual.png}
    \caption{Coeficientes estandarizados del modelo OLS para PD.}
    \label{fig:coef_pd}
\end{figure}

\paragraph{Análisis por variable}

\textbf{VERB\_Ratio} presenta un coeficiente negativo y significativo ($\beta = -0.1329$, $p = 0.020$). Una mayor proporción de verbos se asocia con menor poder discriminativo. Esto sugiere que las pistas centradas en acciones tienden a activar características compartidas entre múltiples candidatos, reduciendo la separación relativa en el espacio de embeddings.

\textbf{Max\_IDF} muestra también un efecto negativo significativo ($\beta = -0.1000$, $p = 0.010$). La mayor rareza léxica no incrementa la discriminación; por el contrario, se asocia con una ligera reducción del PD. Este resultado indica que la especificidad estadística de un término no garantiza su relevancia discriminativa respecto al objetivo.

\textbf{ADJ\_Ratio} ($p = 0.071$) y \textbf{NER\_Count} ($p = 0.068$) presentan efectos positivos marginales, aunque no alcanzan el umbral convencional de significancia del 5\%. Ambos podrían estar relacionados con mayor anclaje descriptivo, pero la evidencia estadística es débil.

\textbf{PROPN\_Ratio}, \textbf{NUM\_Ratio} y \textbf{Tree\_Depth} no muestran asociación significativa con PD en este modelo.

\paragraph{Diagnóstico}

El análisis gráfico (residuos vs ajustados y Q-Q plot) no muestra patrones sistemáticos evidentes. Se observa ligera desviación en colas del Q-Q plot, esperable en métricas acotadas. La utilización de errores robustos HC3 mitiga posibles efectos de heterocedasticidad.

\subsubsection{Ambigüedad (AMB)}

La métrica de Ambigüedad (AMB) cuantifica cuántos candidatos son activados simultáneamente por una pista, ponderando dicha activación por la similitud semántica en el espacio de embeddings. Valores altos indican que la pista distribuye probabilidad sobre múltiples personajes en lugar de concentrarla en el objetivo.

\paragraph{Bondad de ajuste}

El modelo OLS estimado sobre $N = 360$ observaciones produce los siguientes resultados globales:

\begin{itemize}
\item $R^2 = 0.022$
\item $R^2$ ajustado = 0.002
\item $F = 3.042$
\item $p$-valor global = 0.00405
\end{itemize}

El modelo resulta globalmente significativo, aunque su capacidad explicativa es muy reducida (2.2\% de la varianza). Esto indica que las características lingüísticas contienen señal estadística débil respecto a la ambigüedad, pero dicha señal es estructuralmente limitada.

\paragraph{Coeficientes estandarizados}

\begin{table}[H]
\centering
\begin{tabular}{lccc}
\hline
Variable & $\beta$ & $p$-valor & Signif. \\
\hline
PROPN\_Ratio & -0.0284 & 0.649 & -- \\
NUM\_Ratio & -0.0316 & 0.620 & -- \\
Max\_IDF & 0.0755 & 0.002 & * \\
ADJ\_Ratio & -0.0802 & 0.105 & -- \\
VERB\_Ratio & 0.0565 & 0.282 & -- \\
Tree\_Depth & -0.0203 & 0.636 & -- \\
NER\_Count & 0.0462 & 0.539 & -- \\
\hline
\end{tabular}
\caption{Coeficientes estandarizados del modelo OLS para AMB (errores robustos HC3).}
\end{table}

(* $p<0.05$)

La Figura \ref{fig:coef_amb} muestra los coeficientes estandarizados junto con sus intervalos de confianza.

\begin{figure}[H]
    \centering
    \includegraphics[width=0.9\textwidth]{imagenes/ols_AMB_individual.png}
    \caption{Coeficientes estandarizados del modelo OLS para AMB.}
    \label{fig:coef_amb}
\end{figure}

\paragraph{Análisis por variable}

El único predictor estadísticamente significativo es \textbf{Max\_IDF} ($\beta = 0.0755$, $p = 0.002$). La mayor rareza léxica se asocia con un incremento de la ambigüedad. Esto sugiere que términos infrecuentes no necesariamente restringen el campo semántico, sino que pueden activar regiones menos densas pero más dispersas del espacio de embeddings, aumentando el solapamiento relativo entre candidatos.

El resto de variables no presenta efectos significativos:

\begin{itemize}
\item \textbf{PROPN\_Ratio} y \textbf{NUM\_Ratio} no muestran reducción sistemática de ambigüedad.
\item \textbf{VERB\_Ratio} y \textbf{ADJ\_Ratio} no presentan asociación robusta.
\item \textbf{Tree\_Depth} y \textbf{NER\_Count} no muestran efectos estructurales detectables.
\end{itemize}

\paragraph{Diagnóstico}

El análisis gráfico no muestra patrones claros de heterocedasticidad o no linealidad. Se observa ligera desviación en colas del Q-Q plot, consistente con la naturaleza acotada de la métrica. La utilización de errores robustos HC3 garantiza inferencia válida ante posibles desviaciones de homocedasticidad.

En conjunto, los resultados indican que la ambigüedad es escasamente explicable a partir de propiedades sintácticas superficiales. La baja varianza explicada sugiere que su comportamiento depende en mayor medida de la estructura relacional del espacio semántico que de la forma lingüística del enunciado.

\subsubsection{Ganancia de Información (nGAIN)}

La métrica nGAIN cuantifica la reducción de entropía inducida por la pista en la distribución de probabilidad sobre el conjunto de candidatos. Se encuentra normalizada en el intervalo $[0,1]$ y mide cuánto se concentra la probabilidad tras observar la pista respecto a la distribución inicial.

Desde una perspectiva informacional, nGAIN captura la magnitud absoluta del desplazamiento probabilístico, independientemente de cuál sea el competidor más cercano.

\paragraph{Bondad de ajuste}

El modelo OLS estimado sobre $N = 360$ observaciones produce los siguientes resultados globales:

\begin{itemize}
\item $R^2 = 0.061$
\item $R^2$ ajustado = 0.042
\item $F = 3.342$
\item $p$-valor global = 0.00184
\end{itemize}

El modelo resulta globalmente significativo y explica aproximadamente un 6.1\% de la varianza de nGAIN. Aunque la capacidad explicativa es moderada, es superior a la observada en AMB y ligeramente mayor que en PD, lo que sugiere que la ganancia de información es parcialmente modulable desde la estructura lingüística.

\paragraph{Coeficientes estandarizados}

\begin{table}[H]
\centering
\begin{tabular}{lccc}
\hline
Variable & $\beta$ & $p$-valor & Signif. \\
\hline
PROPN\_Ratio & -0.0064 & 0.920 & -- \\
NUM\_Ratio & -0.0335 & 0.646 & -- \\
Max\_IDF & -0.0241 & 0.636 & -- \\
ADJ\_Ratio & 0.0136 & 0.820 & -- \\
VERB\_Ratio & -0.2092 & 0.000 & *** \\
Tree\_Depth & -0.0550 & 0.285 & -- \\
NER\_Count & 0.1240 & 0.029 & * \\
\hline
\end{tabular}
\caption{Coeficientes estandarizados del modelo OLS para nGAIN (errores robustos HC3).}
\end{table}

($*$ $p<0.05$, $***$ $p<0.001$)

La Figura \ref{fig:coef_ngain} muestra los coeficientes estandarizados junto con sus intervalos de confianza.

\begin{figure}[H]
    \centering
    \includegraphics[width=0.9\textwidth]{imagenes/ols_nGAIN_individual.png}
    \caption{Coeficientes estandarizados del modelo OLS para nGAIN.}
    \label{fig:coef_ngain}
\end{figure}

\paragraph{Análisis por variable}

\textbf{VERB\_Ratio} presenta el efecto más fuerte del modelo ($\beta = -0.2092$, $p < 0.001$). Una mayor proporción de verbos se asocia con una reducción significativa de la ganancia de información. Esto indica que las pistas centradas en acciones tienden a distribuir probabilidad sobre múltiples candidatos, disminuyendo la reducción neta de entropía. En términos estructurales, los verbos parecen introducir señal compartida más que discriminativa.

\textbf{NER\_Count} muestra un efecto positivo significativo ($\beta = 0.1240$, $p = 0.029$). La presencia de entidades nombradas incrementa la ganancia de información. Las entidades concretas actúan como anclajes semánticos que restringen el espacio de búsqueda de manera más precisa que formulaciones genéricas.

El resto de variables no presenta asociación estadísticamente significativa:

\begin{itemize}
\item \textbf{PROPN\_Ratio} y \textbf{NUM\_Ratio} no muestran efecto detectable sobre la reducción de entropía.
\item \textbf{Max\_IDF} no evidencia relación sistemática con la ganancia informacional.
\item \textbf{ADJ\_Ratio} no presenta efecto estructural relevante.
\item \textbf{Tree\_Depth} no muestra impacto significativo.
\end{itemize}

\paragraph{Diagnóstico}

Los gráficos de residuos no evidencian patrones sistemáticos de no linealidad. Se observa ligera desviación en colas del Q-Q plot, coherente con métricas normalizadas. La utilización de errores robustos HC3 garantiza inferencia consistente ante posibles desviaciones de homocedasticidad.

En conjunto, los resultados muestran que la reducción efectiva de entropía depende principalmente de la presencia de entidades explícitas y de la moderación en el uso de verbos, mientras que otras dimensiones superficiales de la sintaxis no presentan efectos sistemáticos detectables.

\subsubsection{Lecturabilidad (LEC)}

La métrica LEC se basa en el índice de Flesch Reading Ease adaptado al español y mide la facilidad cognitiva de procesamiento del texto. Valores altos indican mayor simplicidad sintáctica y menor carga estructural.

\paragraph{Bondad de ajuste}

El modelo OLS estimado sobre $N = 360$ observaciones produce los siguientes resultados globales:

\begin{itemize}
\item $R^2 = 0.154$
\item $R^2$ ajustado = 0.137
\item $F = 3.790$
\item $p$-valor global = 0.000556
\end{itemize}

El modelo resulta altamente significativo y explica aproximadamente un 15.4\% de la varianza en LEC. Es, hasta este punto, el modelo con mayor capacidad explicativa entre las métricas analizadas, lo que es coherente con la naturaleza estructural de la lecturabilidad.

\paragraph{Coeficientes estandarizados}

\begin{table}[H]
\centering
\begin{tabular}{lccc}
\hline
Variable & $\beta$ & $p$-valor & Signif. \\
\hline
PROPN\_Ratio & -0.0769 & 0.101 & -- \\
NUM\_Ratio & 0.0379 & 0.190 & -- \\
Max\_IDF & -0.0368 & 0.296 & -- \\
ADJ\_Ratio & -0.3793 & 0.000 & *** \\
VERB\_Ratio & -0.2141 & 0.001 & ** \\
Tree\_Depth & -0.0432 & 0.337 & -- \\
NER\_Count & 0.0617 & 0.160 & -- \\
\hline
\end{tabular}
\caption{Coeficientes estandarizados del modelo OLS para LEC (errores robustos HC3).}
\end{table}

($**$ $p<0.01$, $***$ $p<0.001$)

La Figura \ref{fig:coef_lec} muestra los coeficientes estandarizados junto con sus intervalos de confianza.

\begin{figure}[H]
    \centering
    \includegraphics[width=0.9\textwidth]{imagenes/ols_LEC_individual.png}
    \caption{Coeficientes estandarizados del modelo OLS para LEC.}
    \label{fig:coef_lec}
\end{figure}

\paragraph{Análisis por variable}

\textbf{ADJ\_Ratio} presenta el mayor efecto estadístico del modelo ($\beta = -0.3793$, $p < 0.001$). Un incremento de una desviación estándar en la proporción de adjetivos se asocia empíricamente con una reducción de aproximadamente 0.38 desviaciones estándar en la lecturabilidad. Este es el coeficiente individual de mayor magnitud observado en el conjunto de regresiones estimadas, indicando que una mayor densidad de modificadores atributivos correlaciona fuertemente con una menor facilidad de lectura.

\textbf{VERB\_Ratio} también muestra un efecto negativo y estadísticamente significativo ($\beta = -0.2141$, $p = 0.001$). Una mayor proporción de verbos en la pista se asocia sistemáticamente con una disminución en la métrica de lecturabilidad, constituyendo el segundo efecto de mayor peso dentro de este modelo en particular.

El resto de las variables evaluadas no presentan efectos estadísticamente significativos sobre la lecturabilidad en esta muestra. Concretamente, \textbf{NUM\_Ratio} ($p = 0.190$) y \textbf{Tree\_Depth} ($p = 0.337$) no muestran una asociación sistemática con la variable dependiente, al igual que ocurre con \textbf{PROPN\_Ratio}, \textbf{Max\_IDF} y \textbf{NER\_Count}.

\paragraph{Diagnóstico}

La evaluación formal de los residuos revela desviaciones respecto a los supuestos clásicos de regresión. El test de Shapiro-Wilk ($W = 0.6600$, $p < 0.001$) indica que los residuos no se distribuyen de forma normal. Asimismo, el test de Breusch-Pagan ($LM = 28.3958$, $p < 0.001$) detecta la presencia de heterocedasticidad en el modelo.

Ante este diagnóstico, la decisión metodológica de emplear errores estándar robustos HC3 resulta fundamental. Esta corrección garantiza que la inferencia estadística realizada sobre los coeficientes (específicamente los intervalos de confianza y los $p$-valores asociados a \textbf{ADJ\_Ratio} y \textbf{VERB\_Ratio}) mantenga su validez matemática a pesar de la ausencia de homocedasticidad y normalidad en los residuos.

\subsubsection{Consistencia (CONS)}

La métrica de Consistencia (CONS) mide la coherencia semántica interna del conjunto de pistas asociadas a un personaje, evaluando cuantitativamente si todas ellas convergen hacia un mismo eje temático o si presentan dispersión.

\paragraph{Bondad de ajuste}

El modelo OLS estimado sobre $N = 360$ observaciones produce los siguientes resultados globales:

\begin{itemize}
\item $R^2 = 0.138$
\item $R^2$ ajustado = 0.121
\item $F = 7.285$
\item $p$-valor global $< 0.001$ ($3.63 \times 10^{-8}$)
\end{itemize}

El modelo resulta globalmente significativo y explica aproximadamente un 13.8\% de la varianza en CONS. Esta capacidad explicativa moderada sugiere que la coherencia temática está asociada de forma detectable con ciertas decisiones de estructuración lingüística.

\paragraph{Coeficientes estandarizados}

La siguiente tabla resume los coeficientes estimados utilizando errores robustos HC3:

\begin{table}[H]
\centering
\begin{tabular}{lccc}
\hline
Variable & $\beta$ & $p$-valor & Signif. \\
\hline
PROPN\_Ratio & -0.0904 & 0.230 & -- \\
NUM\_Ratio & -0.0884 & 0.180 & -- \\
Max\_IDF & -0.0290 & 0.270 & -- \\
ADJ\_Ratio & -0.0661 & 0.333 & -- \\
VERB\_Ratio & -0.3016 & 0.000 & *** \\
Tree\_Depth & 0.1131 & 0.020 & * \\
NER\_Count & -0.0977 & 0.147 & -- \\
\hline
\end{tabular}
\caption{Coeficientes estandarizados del modelo OLS para CONS (errores robustos HC3).}
\end{table}

(* $p<0.05$, *** $p<0.001$)

La Figura \ref{fig:coef_cons} muestra los coeficientes estandarizados junto con sus intervalos de confianza.

\begin{figure}[H]
    \centering
    \includegraphics[width=0.9\textwidth]{imagenes/ols_CONS_individual.png}
    \caption{Coeficientes estandarizados del modelo OLS para CONS.}
    \label{fig:coef_cons}
\end{figure}

\paragraph{Análisis por variable}

\textbf{VERB\_Ratio} presenta el efecto negativo más intenso de este modelo ($\beta = -0.3016$, $p < 0.001$). Un incremento en la proporción de verbos se asocia con una reducción significativa de la consistencia semántica entre las pistas. Estadísticamente, la introducción de una mayor densidad verbal correlaciona con una mayor dispersión direccional en el espacio temático.

\textbf{Tree\_Depth} muestra un efecto positivo y estadísticamente significativo ($\beta = 0.1131$, $p = 0.020$). Una mayor profundidad del árbol sintáctico se asocia con un incremento en la consistencia. Este resultado indica que las oraciones con mayor complejidad estructural o anidamiento tienden a presentar una organización temática internamente más cohesionada.

El resto de las variables (\textbf{PROPN\_Ratio}, \textbf{NUM\_Ratio}, \textbf{Max\_IDF}, \textbf{ADJ\_Ratio} y \textbf{NER\_Count}) no presentan asociaciones estadísticamente significativas con la consistencia en esta muestra, manteniendo $p$-valores ampliamente superiores al nivel $\alpha = 0.05$.

\paragraph{Diagnóstico}

La evaluación de los residuos sugiere un comportamiento estadístico manejable dentro de los parámetros del modelo. Como en las estimaciones anteriores, la utilización de errores estándar robustos HC3 garantiza que la inferencia sobre los coeficientes (\textit{p}-valores e intervalos de confianza) sea consistente y válida ante posibles desviaciones en los supuestos de normalidad u homocedasticidad.


\subsection{H4 --- Efecto de la modalidad experimental}

La segunda fase experimental introduce una variación fundamental en la mecánica de juego: los jugadores ya no recuperan el nombre del personaje de memoria, sino que lo seleccionan de un conjunto cerrado de candidatos predefinidos. Este cambio no afecta a las pistas en sí, que siguen siendo los mismos 360 enunciados generados con idéntico procedimiento, pero transforma radicalmente la dinámica cognitiva del acierto. El presente análisis se organiza en torno a esta asimetría: los modelos OLS son matemáticamente invariantes (las métricas de calidad son propiedades del texto, no del jugador), mientras que la regresión logística refleja un paisaje predictivo sustancialmente distinto al observado en mundo abierto.

% -------------------------------------------------------
\subsubsection{Invarianza del análisis OLS}

Las métricas de calidad PD, AMB, nGAIN, LEC y CONS se calculan sobre las pistas como entidades lingüísticas: sus valores dependen de los embeddings y la estructura sintáctica de cada enunciado, con independencia de cómo el jugador interactúe con él. Al ser las 360 pistas idénticas en ambas fases, los cinco modelos OLS producen exactamente los mismos coeficientes, los mismos valores de $R^2$ y los mismos niveles de significancia. Repetir el análisis OLS en esta sección no aportaría ningún resultado nuevo.

Esta invarianza no es un resultado vacío: confirma que la relación entre estructura lingüística y calidad semántica es una propiedad estable del texto. Los coeficientes y su interpretación se encuentran en la subsección correspondiente a la Fase 1 y aplican sin modificación al presente experimento.

% -------------------------------------------------------
\subsubsection{Regresión logística: Fase 2 (Mundo Cerrado)}

\paragraph{Modalidad de respuesta y dinámica del acierto}

En mundo cerrado el jugador dispone de una lista visible de candidatos y debe identificar cuál de ellos corresponde a la pista. Esta condición elimina la necesidad de recuperación libre de memoria y sustituye la inferencia abierta por un proceso de discriminación comparativa: el jugador puede evaluar cada candidato de la lista frente al contenido de la pista y seleccionar el más probable. Esta diferencia cognitiva se traduce directamente en una tasa de acierto mucho más elevada (78.1\% frente al 65.5\% observado en mundo abierto) y en un patrón de importancia de métricas cualitativamente distinto.

\paragraph{Datos y distribución de clases}

El conjunto de datos de esta fase comprende \textbf{$N = 1\,119$} registros de partida, recogidos bajo la modalidad de selección restringida. La distribución de clases es la siguiente:

\begin{table}[H]
\centering
\begin{tabular}{lcc}
\hline
Clase & Frecuencia & Porcentaje \\
\hline
Acierto (1) & 874 & 78.1\% \\
Fallo (0)   & 245 & 21.9\% \\
\hline
Total & 1\,119 & \\
\hline
\end{tabular}
\caption{Distribución de clases en el conjunto de datos de Fase 2.}
\label{tab:clases_f2}
\end{table}

El desbalance de clases (casi 4:1 a favor del acierto) es más pronunciado que en mundo abierto (aproximadamente 1.9:1). Esto impone restricciones a la capacidad del modelo para detectar los casos de fallo, que constituyen la señal más informativa para la predicción. El parámetro \texttt{class\_weight='balanced'} mitiga parcialmente este efecto pero no lo elimina.

\paragraph{Bondad de ajuste y validación cruzada}

El modelo se evalúa mediante validación cruzada estratificada de 5 pliegues y un conjunto de prueba retenido (80/20 estratificado). La Tabla~\ref{tab:bondad_f2} compara las métricas de ajuste con las obtenidas en mundo abierto:

\begin{table}[H]
\centering
\begin{tabular}{lccc}
\hline
Métrica & Mundo Cerrado (F2) & Mundo Abierto (F1) & Diferencia \\
\hline
Accuracy (CV 5-fold)   & $0.532 \pm 0.040$ & $\sim 0.55 \pm 0.03$ & $-0.018$ \\
AUC-ROC (CV 5-fold)    & $0.548 \pm 0.058$ & $\sim 0.565 \pm 0.027$ & $-0.017$ \\
AUC-ROC (test set)     & $0.505$           & $0.591$              & $-0.086$ \\
McFadden $R^2$         & $0.0132$          & $0.0102$             & $+0.003$ \\
McFadden $R^2$ adj.    & $0.0030$          & $0.0061$             & $-0.003$ \\
LR test $G^2$          & $15.56$           & $30.10$              & $-14.54$ \\
LR test $p$            & $0.0082$          & $< 0.001$            & — \\
Brier Score            & $0.2495$          & $0.245$              & $+0.005$ \\
\hline
\end{tabular}
\caption{Comparación de bondad de ajuste del modelo logístico entre ambas fases.}
\label{tab:bondad_f2}
\end{table}

A pesar de que los jugadores aciertan con más frecuencia en mundo cerrado, el modelo tiene más dificultades para predecir el éxito con las métricas de la pista. La paradoja se explica por el extremo desbalance de clases: cuando casi todos los jugadores aciertan, el margen de variación atribuible a la calidad de la pista es mucho menor. El modelo logístico sigue siendo estadísticamente significativo ($G^2 = 15.56$, $p = 0.0082$), lo que confirma que las métricas de calidad contienen señal predictiva real incluso en este contexto más restrictivo.

\begin{figure}[H]
    \centering
    \includegraphics[width=\textwidth]{imagenes/logistic_validation_completa.png}
    \caption{Validación completa del modelo logístico en Fase 2. Se incluyen la curva ROC sobre el conjunto de prueba, las métricas de validación cruzada estratificada (accuracy, AUC, F1, log-loss) y el reporte de clasificación. La curva ROC muestra capacidad discriminativa modesta pero consistentemente superior al azar.}
    \label{fig:logistic_val_f2}
\end{figure}

\begin{figure}[H]
    \centering
    \includegraphics[width=0.75\textwidth]{imagenes/roc_curve_cv_inc.png}
    \caption{Curva ROC con validación cruzada de 5 pliegues (Fase 2). La banda sombreada representa la desviación estándar entre pliegues. El AUC medio de $0.548 \pm 0.058$ es inferior al obtenido en mundo abierto ($0.565 \pm 0.027$) y presenta mayor varianza entre pliegues, reflejo del menor tamaño muestral.}
    \label{fig:roc_cv_f2}
\end{figure}

\paragraph{Coeficientes e importancia de métricas}

La Tabla~\ref{tab:coef_logistica_f2} presenta los coeficientes $\beta$ estandarizados del modelo sklearn (con regularización L2) y los contrasta con los de mundo abierto, revelando las diferencias en la estructura de importancia entre modalidades:

\begin{table}[H]
\centering
\begin{tabular}{lcccc}
\hline
Métrica & $\beta$ F2 & Dirección & $\beta$ F1 (ref.) & Cambio \\
\hline
\textbf{PD}   & $\mathbf{+0.333}$ & $\uparrow$ acierto & $+0.110$ & $\times 3.0$  \\
\textbf{CONS} & $-0.162$          & $\downarrow$ acierto & $-0.008$ & emerge \\
\textbf{nGAIN}& $+0.061$          & $\uparrow$ acierto & $+0.056$ & estable \\
\textbf{AMB}  & $+0.051$          & $\uparrow$ acierto & $-0.080$ & \textit{invierte} \\
\textbf{LEC}  & $+0.023$          & $\uparrow$ acierto & $-0.070$ & \textit{invierte} \\
\hline
\end{tabular}
\caption{Coeficientes $\beta$ estandarizados del modelo logístico (sklearn, L2) en Fase 2, comparados con los de Fase 1. Los pesos normalizados derivados sitúan PD en el 53.2\%, CONS en el 25.9\%, nGAIN en el 9.8\%, AMB en el 8.1\% y LEC en el 3.7\%.}
\label{tab:coef_logistica_f2}
\end{table}

En cuanto a la significancia individual (modelo statsmodels, sin regularización), únicamente \textbf{PD} alcanza el umbral convencional del 5\% ($\beta = +0.325$, $p = 0.028$). El resto de coeficientes no son individualmente significativos, en parte por el menor tamaño muestral ($N = 1\,119$ frente a $N = 2\,293$), aunque el modelo conjunto sigue siendo significativo ($p = 0.008$).

\begin{figure}[H]
    \centering
    \includegraphics[width=0.8\textwidth]{imagenes/logistic_importance.png}
    \caption{Importancia relativa de las métricas de calidad en el modelo logístico de Fase 2 (pesos normalizados a partir de $|\beta|$). PD absorbe más del 53\% del peso total, mientras que en Fase 1 su contribución era de aproximadamente el 34\%. CONS emerge como segundo factor relevante (25.9\%) desde una posición prácticamente irrelevante en mundo abierto (2\%).}
    \label{fig:logistic_importance_f2}
\end{figure}

\paragraph{Análisis por métrica}

\textbf{PD: de relevante a dominante.} En mundo abierto, el Poder Discriminativo tenía el coeficiente positivo más alto pero no se distinguía estadísticamente del ruido. En mundo cerrado, PD triplica su coeficiente ($+0.110 \rightarrow +0.333$) y se convierte en el único predictor individualmente significativo. La explicación es directa: en presencia de una lista de candidatos el jugador realiza explícitamente un proceso de discriminación comparativa: evalúa cada candidato frente a la pista, lo que hace que la capacidad de la pista para separar el target del resto de candidatos sea la propiedad más directamente útil. Una pista con PD alto apunta de forma nítida hacia un individuo concreto dentro de la lista visible; una pista con PD bajo distribuye similitud entre varios candidatos y dificulta la elección.

\textbf{AMB: inversión de signo por mediación de la lista.} En mundo abierto, la ambigüedad semántica tenía un efecto negativo sobre el acierto ($\beta = -0.080$): cuando el jugador tiene que recuperar el nombre de memoria, una pista que activa múltiples candidatos lo lleva a callejones sin salida. En mundo cerrado el signo se invierte ($\beta = +0.051$): la lista de candidatos actúa como filtro externo que resuelve la ambigüedad. Aunque la pista no apunte limpiamente hacia el target, el jugador puede confrontarla contra cada candidato de la lista y descartar los que son claramente incompatibles. La ambigüedad intrínseca de la pista importa mucho menos cuando el proceso de eliminación lo guía la interfaz, no la memoria.

\textbf{LEC: dilución del efecto en mundo cerrado.} La Lecturabilidad presentaba un coeficiente negativo en mundo abierto ($\beta = -0.070$), interpretado como señal de que el vocabulario técnico y específico, que penaliza el índice Flesch, era en realidad más discriminativo. En mundo cerrado ese efecto desaparece casi por completo ($\beta = +0.023$, no significativo). Cuando el jugador puede leer la pista tranquilamente y compararla con una lista finita, la carga cognitiva del vocabulario difícil deja de ser un obstáculo: tiene tiempo para procesar cada término y confrontarlo con los candidatos disponibles. La ventaja del vocabulario técnico se diluye.

\textbf{CONS: de irrelevante a segundo factor.} La Consistencia interna tenía un coeficiente prácticamente nulo en mundo abierto ($\beta = -0.008$). En mundo cerrado emerge como segundo predictor en magnitud ($\beta = -0.162$). Una hipótesis explicativa es que cuando las pistas de un personaje son muy temáticamente cohesionadas, pueden crear un perfil genérico dentro de su categoría. Si todos los clues de un personaje hablan de ``actuación y premios cinematográficos'', ese perfil puede describir con igual plausibilidad a varios candidatos del panel que pertenecen a la misma categoría. La consistencia alta no ayuda a diferenciar al individuo concreto de sus vecinos semánticos más próximos en la lista.

\textbf{nGAIN: el único coeficiente estable.} La Ganancia de Información Normalizada es la única métrica que no cambia sustancialmente entre fases ($+0.056 \rightarrow +0.061$). Esto sugiere que la capacidad de la pista para reducir la entropía del espacio de candidatos es igualmente valiosa independientemente de si el jugador accede a ese espacio de forma libre o restringida. Una pista que concentra probabilidad sobre el target lo hace tanto si el jugador recuerda el nombre por sí solo como si lo selecciona de una lista.

% -------------------------------------------------------
\subsubsection{Comparación de las tres variantes de $Q$}

La función de calidad compuesta $Q$ existe en tres versiones: la \textbf{$Q$ Inicial}, calculada con los pesos teóricos asignados antes de cualquier análisis empírico ($w_{nGAIN}=0.30$, $w_{PD}=w_{AMB}=0.20$, $w_{CONS}=w_{LEC}=0.15$); la \textbf{$Q$ Mundo Abierto}, con pesos derivados de la regresión logística de Fase 1 ($w_{PD}=0.335$, $w_{LEC}=0.303$, $w_{AMB}=0.215$, $w_{nGAIN}=0.139$, $w_{CONS}=0.009$); y la \textbf{$Q$ Mundo Cerrado}, con los pesos de Fase 2 ($w_{PD}=0.529$, $w_{CONS}=0.257$, $w_{nGAIN}=0.097$, $w_{AMB}=0.081$, $w_{LEC}=0.037$). Las figuras siguientes comparan las tres distribuciones sobre las mismas 360 pistas.

\begin{figure}[H]
    \centering
    \includegraphics[width=\textwidth]{imagenes/q_comparison_inc.png}
    \caption{Panel resumen ($2\times3$) con las seis comparaciones entre las tres variantes de $Q$. Fila superior: histogramas superpuestos, box plots y violin plots. Fila inferior: diagramas de dispersión para cada par de variantes ($Q$~Inicial vs $Q$~M.Abierto, $Q$~Inicial vs $Q$~M.Cerrado, $Q$~M.Abierto vs $Q$~M.Cerrado). Azul: $Q$ Inicial; morado: $Q$ Mundo Abierto; verde: $Q$ Mundo Cerrado.}
    \label{fig:q_comparison_f2}
\end{figure}

\begin{figure}[H]
    \centering
    \includegraphics[width=\textwidth]{imagenes/q_comparison_histogramas_inc.png}
    \caption{Histogramas superpuestos de las tres variantes de $Q$ sobre las 360 pistas. La $Q$ Inicial (azul) concentra masa en una banda estrecha centrada en valores medios, reflejo de que todos los pesos eran moderados. La $Q$ Mundo Abierto (morado) se desplaza según la doble dominancia de PD y LEC. La $Q$ Mundo Cerrado (verde) muestra la mayor dispersión, consecuencia del dominio extremo de PD (53\%): las pistas con PD alto se separan claramente del resto.}
    \label{fig:q_hist_f2}
\end{figure}

\begin{figure}[H]
    \centering
    \includegraphics[width=0.85\textwidth]{imagenes/q_comparison_boxplot_inc.png}
    \caption{Diagramas de caja comparativos de las tres variantes de $Q$. La comparación de medianas y rangos intercuartílicos revela cómo el rebalanceo de pesos entre variantes desplaza la tendencia central y modifica la asimetría de la distribución. La mayor varianza de $Q$ Mundo Cerrado es consistente con la polarización inducida por el peso dominante de PD.}
    \label{fig:q_boxplot_f2}
\end{figure}

\begin{figure}[H]
    \centering
    \includegraphics[width=0.85\textwidth]{imagenes/q_comparison_violin_inc.png}
    \caption{Diagramas de violín de las tres variantes de $Q$. La forma del violín revela la densidad de probabilidad en cada punto del rango de valores. La $Q$ Inicial muestra una forma más compacta y simétrica; la $Q$ Mundo Cerrado presenta una forma más elongada con un pico pronunciado en valores medios-bajos y cola más larga en los valores altos, asociada a las pistas con PD excepcional.}
    \label{fig:q_violin_f2}
\end{figure}

\begin{figure}[H]
    \centering
    \includegraphics[width=\textwidth]{imagenes/q_comparison_scatter_inc.png}
    \caption{Diagramas de dispersión para cada par de variantes. \textit{Izquierda}: $Q$ Inicial vs $Q$ Mundo Abierto — la correlación es moderada; las pistas con LEC alta ganan posición respecto a la $Q$ Inicial. \textit{Centro}: $Q$ Inicial vs $Q$ Mundo Cerrado — la dispersión alrededor de la diagonal es mayor, indicando un reordenamiento más profundo inducido por el peso dominante de PD. \textit{Derecha}: $Q$ Mundo Abierto vs $Q$ Mundo Cerrado — el tercer scatter muestra directamente cuánto difiere el ordenamiento de las pistas entre las dos fórmulas empíricas: las pistas que tienen PD alto pero LEC baja suben en mundo cerrado y bajan en mundo abierto, y viceversa.}
    \label{fig:q_scatter_f2}
\end{figure}

\begin{table}[H]
\centering
\begin{tabular}{lccc}
\hline
Estadístico & $Q$ Inicial & $Q$ M.\,Abierto & $Q$ M.\,Cerrado \\
\hline
Media         & 0.344     & 0.292     & 0.277  \\
Mediana       & 0.345     & 0.301     & 0.282  \\
Desv.\,típica & 0.060     & 0.070     & 0.052  \\
Mínimo        & $-$0.068  & $-$0.232  & 0.027  \\
Máximo        & 0.507     & 0.466     & 0.404  \\
$P_{25}$      & 0.315     & 0.261     & 0.247  \\
$P_{75}$      & 0.378     & 0.331     & 0.310  \\
\hline
\multicolumn{4}{l}{\textit{Correlaciones de Pearson entre variantes}} \\
\hline
$Q$ Inicial $\leftrightarrow$ $Q$ M.\,Abierto     & \multicolumn{3}{c}{$r = 0.919$} \\
$Q$ Inicial $\leftrightarrow$ $Q$ M.\,Cerrado     & \multicolumn{3}{c}{$r = 0.884$} \\
$Q$ M.\,Abierto $\leftrightarrow$ $Q$ M.\,Cerrado & \multicolumn{3}{c}{$r = 0.937$} \\
\hline
\end{tabular}
\caption{Estadísticos descriptivos de las tres variantes de $Q$ sobre las 360 pistas del corpus, con correlaciones de Pearson entre pares ($N=360$).}
\label{tab:q_stats_f2}
\end{table}

Las tres correlaciones son elevadas ($r \geq 0.88$), lo que indica que el rebalanceo de pesos no altera radicalmente el ordenamiento de las pistas: las pistas valoradas como buenas por la fórmula inicial tienden a serlo también por ambas fórmulas empíricas. La correlación más baja ($r = 0.884$ entre $Q$ Inicial y $Q$ M.\,Cerrado) refleja que el cambio de paradigma modifica de forma más pronunciada la valoración relativa de las pistas que la transición entre dos configuraciones de mundo abierto. Sin embargo, los estadísticos revelan desplazamientos sistemáticos: las fórmulas empíricas reducen la media en conjunto (de 0.344 a 0.292 y 0.277), reflejo de que la $Q$ Inicial era más generosa al distribuir el peso en nGAIN, métrica con valores típicamente elevados. La $Q$ Mundo Cerrado presenta la menor dispersión ($\sigma = 0.052$) y el mínimo absoluto más alto (0.027, sin valores negativos), efecto directo de que todos sus pesos apuntan en el mismo sentido en la fórmula compuesta: PD y nGAIN en positivo, CONS y AMB como penalizaciones suaves vía el término $(1-1.5 \cdot \text{AMB})$.

La evolución de las tres variantes ilustra el impacto empírico de los datos de juego sobre la fórmula de calidad. Mientras que la $Q$ Inicial distribuía el peso de forma relativamente uniforme basándose en criterios teóricos, ambas versiones empíricas convergen en elevar el papel de PD, aunque en distinta medida, y reducen drásticamente el peso de nGAIN, que ocupaba el lugar privilegiado en la fórmula original. Las diferencias entre la $Q$ Mundo Abierto y la $Q$ Mundo Cerrado reflejan, pista a pista, la asimetría de las mecánicas: las pistas que el modelo de mundo cerrado valora por encima del modelo de mundo abierto son aquellas que apuntan con precisión al target dentro de una lista, aunque resulten difíciles de procesar o temáticamente dispersas; las pistas que el modelo de mundo abierto prima son aquellas que combinan discriminación semántica con claridad expresiva, necesaria cuando el jugador no dispone de ningún andamiaje externo.
\end{document}